%=======================================================================
%  CHAPTER 5 - FIR FILTER DESIGN  (6 hrs)
%  Every window value, filter length, beta and Bessel figure quoted here
%  is produced by fir.py, which self-tests before it prints.
%  Sources: Chapterwise Ch5 pp.1-27 (windowing pp.1-11, Kaiser pp.12-20
%           with the design equations on p.17, Parks-McClellan pp.21-27
%           with the 6-step procedure and the M=26 vs M=38 comparison
%           on p.27; pp.28-31 are chapter 4 material, not this one);
%           SPP Sir deck 6 "IIR _ FIR" p.123 for the window comparison
%           table (an embedded image of Oppenheim Table 7.2) and p.132
%           for the Kaiser equations, agreeing verbatim with Chapterwise;
%           SPP Sir "Optimum_FIR_Filters" (17 pp, exported this session)
%           for the alternation theorem and the Remez flowchart on p.17.
%           NOTE: DSAP All Chapter Summary p.94 says of Remez only
%           "read yourself from Proakis book", so it is no use here.
%  Numericals live in ch5-num.tex.
%=======================================================================
\section{FIR Filter Design}

{\color{sub}\footnotesize 6 hrs \gap$\cdot$\gap \mk{70 questions across the 45 papers}
\gap$\cdot$\gap worth about \mk{12 of 80 marks} \gap$\cdot$\gap 3 syllabus sub-topics
\gap$\cdot$\gap \mk{4th of the 8 chapters} by weight.}

\lead{Read this first:}
\begin{itemize}
  \item Three topics, and they divide cleanly by \emph{what you have to do}:
  \begin{itemize}
    \item \textbf{Window method} \tS{28}: a design calculation. Pick a window, get a
          length, truncate the ideal response.
    \item \textbf{Kaiser window} \tS{19}: the same calculation with two formulas
          bolted on, for $\beta$ and for the length.
    \item \textbf{Remez exchange} \tS{23}: \mk{pure bookwork}. 23 questions, almost
          always the same flowchart and the same three sentences, and \mk{not one of
          them asks you to run the algorithm}.
  \end{itemize}
  \item So a third of this chapter's marks are free if you memorise one flowchart.
        Do \S5.6 first.
  \item \mk{The whole design method is two lines}: the ideal filter has an infinite,
        non-causal impulse response, so \textbf{truncate it} and \textbf{taper the
        truncation} so the edges do not ring.
  \item The window is chosen by \textbf{stopband attenuation} and the length by
        \textbf{transition width}. Those are two independent lookups and they are the
        only two numbers a fixed-window question turns on.
  \item \mk{Gibbs ripple does not shrink when you lengthen the filter.} That single
        fact is worth 3 to 6 marks in eight papers, and it is the reason Kaiser exists.
\end{itemize}

\hr

\T{5.1 \quad The window method}

\Q{Describe how a digital FIR filter is designed by the window method, and discuss
the characteristics of the different window functions}{\m{4+4}\m{5+3}\m{4}\m{5}}
{\color{sub}\footnotesize \tS{28} \yr{\textbf{70 Ch} \m{4+4}, \textbf{72 Ch} \m{5+3},
\texttt{72 Ma} \m{4}, \textbf{\texttt{70 Bh}} \m{5}}
\gap$\cdot$\gap \mk{the answer is the four steps below, then the table in \S5.2}}

\sbsr{0.5}{%
\lead{Where the ideal response comes from}
An ideal low-pass filter is a brick wall in frequency. Inverse DTFT it:
\[ h_d[n]=\frac{1}{2\pi}\int_{-\omega_c}^{\omega_c}e^{j\omega n}\,d\omega
        =\frac{\sin\omega_c n}{\pi n} \]
\begin{itemize}
  \item This is a $\operatorname{sinc}$: \textbf{infinitely long} and
        \textbf{non-causal}, so it cannot be built.
  \item Delay it by $\alpha=\dfrac{N-1}{2}$ samples to make it causal, which is also
        exactly what gives \textbf{linear phase}:
  \item[] \[ h_d[n]=\frac{\sin\big(\omega_c(n-\alpha)\big)}{\pi(n-\alpha)},\qquad
             h_d[\alpha]=\frac{\omega_c}{\pi} \]
  \item The centre tap is a $0/0$ limit, so \mk{evaluate it separately as
        $\omega_c/\pi$}. Forgetting that is the most common slip in the chapter.
\end{itemize}}{%
\lead{The four steps, which are the whole method}
\begin{enumerate}
  \item Write $h_d[n]$ for the ideal filter, delayed by $\alpha=(N-1)/2$.
  \item \textbf{Choose the window} from the required stopband attenuation.
  \item \textbf{Choose $N$} from the required transition width.
  \item Multiply: $\;h[n]=h_d[n]\,w[n]$ for $0\le n\le N-1$, zero elsewhere.
\end{enumerate}
\lead{Other ideal responses}
\begin{itemize}
  \item \textbf{High pass}: $\delta[n-\alpha]-h_{d,\mathrm{LP}}[n]$, an all-pass minus
        the low pass.
  \item \textbf{Band pass}: $h_{d}(\omega_2)-h_{d}(\omega_1)$, the wider low pass minus
        the narrower.
  \item \mk{Every one of them is symmetric about $\alpha$}, so every windowed FIR
        filter here has linear phase automatically.
\end{itemize}}

\lead{Why multiplying by a window is convolving in frequency}
\begin{itemize}
  \item $h[n]=h_d[n]w[n]$ becomes, in frequency,
  \item[] \[ H(e^{j\omega})=\frac{1}{2\pi}\int_{-\pi}^{\pi}
             H_d(e^{j\theta})\,W(e^{j(\omega-\theta)})\,d\theta \]
  \item So the brick wall is \textbf{smeared by the window's spectrum}. The result is
        the sharp edge replaced by a transition, plus ripple from the window's
        sidelobes.
  \item \mk{Main lobe width sets the transition; sidelobe height sets the ripple.}
        That one sentence explains the entire comparison table in \S5.2.
\end{itemize}

\hr

\T{5.2 \quad The windows, and the Gibbs phenomenon}

\Q{Explain the Gibbs phenomenon: how it arises when the rectangular window is used,
and how it can be minimized}{\m{3+3}\m{5}\m{6}\m{1+4}\m{2+5}}
{\color{sub}\footnotesize \tS{8} \yr{81 Ba, 80 Ba, 73 Shr, 71 Shr,
\textbf{\texttt{79 Ch}}, \textbf{\texttt{74 Bh}}, \textbf{\texttt{71 Bh}},
\textbf{\texttt{81 Ch}}} \gap$\cdot$\gap \mk{asked in eight papers, always the same
answer}}

\sbsr{0.46}{%
\lead{What the Gibbs phenomenon is}
\begin{itemize}
  \item Truncating the sinc with a \textbf{rectangular} window convolves $H_d$ with the
        Dirichlet kernel, whose sidelobes are only \textbf{13 dB} below its main lobe.
  \item The result overshoots the band edge by about \mk{9 \% of the discontinuity
        height}, with ringing either side.
  \item \mk{Increasing $N$ does not help.} The ripples get narrower and more numerous,
        but the overshoot stays at 9 \%. The series converges in the mean-square sense,
        \textbf{not uniformly}.
  \item That is the whole phenomenon, and the fact examiners want stated.
\end{itemize}
\lead{How to minimize it}
\begin{itemize}
  \item \textbf{Use a tapered window.} Hann, Hamming and Blackman fall smoothly to zero
        at the ends instead of being cut off, so their sidelobes are far lower and the
        ripple is far smaller.
  \item The price is a \textbf{wider main lobe}, hence a wider transition band, hence a
        longer filter for the same sharpness.
  \item \mk{Every window is a trade of ripple against transition width.}
\end{itemize}}{%
\lead{The comparison table \mk{(learn this: it answers half the chapter)}}
\par
\begin{tabular}{@{}L{2.0cm}L{1.4cm}L{1.3cm}L{1.3cm}L{1.5cm}@{}}
\toprule
\hrow \thead{Window} & \thead{Main lobe} & \thead{Peak side lobe} &
\thead{Stop band} & \thead{Transition} \\
Rectangular & $\frac{4\pi}{M+1}$ & $-13$ dB & $21$ dB & $1.8\pi/M$ \\
Bartlett & $8\pi/M$ & $-25$ dB & $25$ dB & $6.1\pi/M$ \\
Hann & $8\pi/M$ & $-31$ dB & $44$ dB & $6.2\pi/M$ \\
Hamming & $8\pi/M$ & $-41$ dB & $53$ dB & $6.6\pi/M$ \\
Blackman & $12\pi/M$ & $-57$ dB & $74$ dB & $11\pi/M$ \\
\bottomrule
\end{tabular}
\par\smallskip
\begin{itemize}
  \item \mk{Stopband attenuation is a property of the window alone.} It does not depend
        on $N$ at all, which is exactly why Gibbs ripple cannot be lengthened away.
  \item $M=N-1$ throughout. Papers sometimes call $M$ the length; read which.
  \item The two columns you actually use are \textbf{Stop band} (to choose the window)
        and \textbf{Transition} (to choose $N$).
\end{itemize}}

\lead{The window formulas, for $0\le n\le M$ with $M=N-1$}
\begin{itemize}
  \item[] \[ w_{\text{rect}}[n]=1,\qquad
             w_{\text{Bart}}[n]=1-\frac{|2n-M|}{M} \]
  \item[] \[ w_{\text{Hann}}[n]=0.5-0.5\cos\frac{2\pi n}{M},\qquad
             w_{\text{Hamm}}[n]=0.54-0.46\cos\frac{2\pi n}{M} \]
  \item[] \[ w_{\text{Black}}[n]=0.42-0.5\cos\frac{2\pi n}{M}
             +0.08\cos\frac{4\pi n}{M} \]
  \item \mk{Hann ends at exactly 0; Hamming ends at 0.08.} That is the only difference
        between them and it buys 9 dB of stopband.
  \item All five are \textbf{symmetric}, which is what keeps the phase linear.
\end{itemize}

\hr

\T{5.3 \quad Choosing the window and the length}

\sbs{%
\lead{The two lookups}
\begin{enumerate}
  \item \textbf{Window} $\leftarrow$ stopband attenuation. Take the \emph{cheapest}
        window whose figure meets the spec. If the paper gives a stopband ripple
        $\delta_s$ instead, convert first:
  \item[] \[ A_s=-20\log_{10}\delta_s \ \text{dB} \]
  \item \textbf{Length} $\leftarrow$ transition width
        $\Delta\omega=|\omega_s-\omega_p|$, from the same table:
  \item[] \[ N-1=M=\frac{c\,\pi}{\Delta\omega} \]
        with $c$ the window's transition constant ($6.6$ for Hamming, and so on).
  \item Round $N$ \textbf{up}, and up again to the next \textbf{odd} number so
        $\alpha=(N-1)/2$ is a whole sample.
\end{enumerate}}{%
\lead{The cut-off goes in the middle}
\begin{itemize}
  \item The window method has no independent control of the band edges, so put the
        ideal cut-off \textbf{halfway across the transition band}:
  \item[] \[ \omega_c=\frac{\omega_p+\omega_s}{2} \]
  \item \mk{Use $\omega_c$, not $\omega_p$, in $h_d[n]$.} Using the passband edge is a
        guaranteed lost mark.
  \item If the spec is in Hz, convert first: $\omega=2\pi f/f_s$. If it is in rad/s,
        $\omega=\Omega/f_s$ (equivalently $\Omega T$).
\end{itemize}
{\color{sub}\footnotesize \mk{The window choice is a step, not a slope.} 67 Mng asks
for 54 dB. Hamming gives 53, one short, so Blackman is forced and the filter jumps
from 67 taps to 111 for the sake of 1 dB. Say so: it is the point of the question.}}

\hr

\T{5.4 \quad The Kaiser window}

\Q{Why is the Kaiser window better than other fixed windows? Explain the steps to
design an FIR filter using it}{\m{3}\m{4}\m{7}\m{2+2}}
{\color{sub}\footnotesize \tS{19} \yr{81 Ba, 74 Ash, \textbf{72 Ch},
\textbf{\texttt{80 Ch}}, \textbf{\texttt{72 Ash}}, \textbf{\texttt{74 Bh}} \m{2+2}}
\gap$\cdot$\gap the "why" is one sentence, the design is four lines}

\sbsr{0.46}{%
\lead{Why it is better \mk{(the whole answer)}}
\begin{itemize}
  \item Every fixed window offers \textbf{one} stopband attenuation, take it or leave
        it. Hann gives 44 dB whether you want 30 or 43.
  \item Kaiser has a \textbf{shape parameter $\beta$}, so the attenuation is
        \emph{continuously adjustable}.
  \item So the two specifications \textbf{decouple}: \mk{$\beta$ buys attenuation,
        $N$ buys transition width}, and neither is wasted on the other.
  \item The result is \textbf{the shortest filter} that meets a given spec among all
        the window methods, because nothing is over-designed.
  \item It is also very nearly optimal in the energy sense: it maximises the fraction
        of the window's energy in the main lobe.
\end{itemize}}{%
\lead{The design equations}
\[ w[n]=\frac{I_0\!\left(\beta\sqrt{1-\left(\frac{2n}{M}-1\right)^{2}}\right)}
              {I_0(\beta)},\qquad 0\le n\le M \]
\begin{enumerate}
  \item $\delta=\min(\delta_p,\delta_s)$, then $A=-20\log_{10}\delta$ dB.
  \item[] \[ \beta=\begin{cases}
        0 & A\le21\\
        0.5842(A-21)^{0.4}+0.07886(A-21) & 21<A\le50\\
        0.1102(A-8.7) & A>50\end{cases} \]
  \item[] \[ M=\frac{A-8}{2.285\,\Delta\omega},\qquad \text{length}=M+1 \]
  \item $I_0$ is the \textbf{modified} Bessel function of the first kind, order zero:
        $I_0(x)=\sum_{k\ge0}\left[\frac{(x/2)^k}{k!}\right]^{2}$, and $I_0(0)=1$.
\end{enumerate}
{\color{sub}\footnotesize $\beta=0$ makes the Kaiser window rectangular, and larger
$\beta$ tapers it harder. \mk{The same $\delta$ gives the same $\beta$}, so most papers
in \S2 of the companion share one value.}}

\hr

\T{5.5 \quad Optimum filters and the Remez exchange algorithm}

\Q{What is an optimum filter? Explain the Remez exchange algorithm for FIR filter
design, with its flowchart. Why is it superior to windowing?}{\m{1+5}\m{2+6}\m{4}\m{9}\m{3+10}}
{\color{sub}\footnotesize \tS{23} \yr{\textbf{75 Ch}, \textbf{72 Ch}, \textbf{70 Ch},
\textbf{69 Ch}, \textbf{73 Ch}, \textbf{78 Bh}, \textbf{76 Ch}, \textbf{82 Bh},
\textbf{81 Bh}, 76 Ash, 73 Shr, 71 Shr, \textbf{79 Bh}, 79 Ba, 82 Ba,
\textbf{\texttt{77 Ch}}, \textbf{\texttt{74 Bh}}, \textbf{\texttt{78 Ch}},
\textbf{\texttt{80 Ch}}, \texttt{74 Ma}, \texttt{72 Ma}, \textbf{\texttt{70 Bh}},
\textbf{\texttt{69 Bh}}}
\gap$\cdot$\gap \mk{23 questions, 175 marks, and no calculation in any of them}}

\sbsr{0.48}{%
\lead{What "optimum" means}
\begin{itemize}
  \item The window method accepts whatever error the window happens to give, and that
        error is \textbf{largest at the band edges} and small in the middle. Most of the
        allowed tolerance goes unused.
  \item An \textbf{optimum} (equiripple, minimax, Parks-McClellan) filter instead
        spreads the error evenly, so the worst case is as small as possible:
  \item[] \[ \min_{h}\ \max_{\omega\in S}\ |E(\omega)|,\qquad
             E(\omega)=W(\omega)\big[H_d(\omega)-H(\omega)\big] \]
  \item $W(\omega)$ is a \textbf{weight}, which is how passband and stopband ripple can
        be traded: set $W=\delta_s/\delta_p$ in the passband and $1$ in the stopband.
  \item $S$ is the bands of interest only, so \mk{the transition band is simply left
        out} rather than being forced to behave.
\end{itemize}}{%
\lead{The alternation theorem, which is why it works}
\begin{itemize}
  \item $H$ is the unique best approximation \textbf{if and only if} $E(\omega)$ has at
        least $\mathbf{r+2}$ \textbf{alternations}: extrema of equal magnitude and
        alternating sign, where $r$ is the number of free coefficients.
  \item So the answer is recognisable when you see it: \mk{count the alternations}.
  \item That is what makes the filter \emph{equiripple}: every ripple in a band reaches
        exactly the same height.
\end{itemize}
\lead{Why it beats windowing}
\begin{itemize}
  \item \textbf{Shorter filter} for the same specification. \mk{The local worked
        comparison is worth quoting}: for $\omega_p=0.4\pi$, $\omega_s=0.6\pi$ and
        $\delta=0.0116$, the equiripple design needs $M=26$ while a Kaiser window at
        the same $\delta$ needs $M=38$. \textbf{A third fewer taps.}
  \item \textbf{Independent control} of passband and stopband ripple through $W$.
  \item Handles arbitrary multi-band shapes, not just the standard four.
\end{itemize}}

\lead{The Remez exchange algorithm, as a flowchart \mk{(this is the marks)}}
\begin{enumerate}
  \item \textbf{Input} the band edges, the desired response $H_d$, the weights $W$, and
        the order.
  \item \textbf{Guess} an initial set of $r+2$ extremal frequencies, usually spaced
        evenly over the bands.
  \item \textbf{Solve} for the deviation $\delta$ on that set. It comes out in closed
        form, with no matrix inversion:
  \item[] \[ \delta=\frac{\sum_{k}\gamma_k H_d(\omega_k)}
                       {\sum_{k}(-1)^{k}\gamma_k/W(\omega_k)} \]
  \item \textbf{Interpolate} the trial response $H$ through the points using Lagrange
        interpolation, and evaluate $E(\omega)$ on a dense grid.
  \item \textbf{Exchange}: find the new local extrema of $|E|$ and make them the new
        extremal set.
  \item \textbf{Test}: if the largest $|E|$ on the grid no longer exceeds $|\delta|$,
        the alternation theorem is satisfied, so \textbf{stop}. Otherwise go back to
        step 3.
  \item \textbf{Output} the impulse response from the converged frequency response.
\end{enumerate}
{\color{sub}\footnotesize Each pass can only increase $|\delta|$ towards its optimum, so
the iteration always converges. \mk{Draw it as a loop with the test at the bottom},
label the exchange step, and the flowchart marks are safe. No paper in the 45 asks for
a numerical iteration.}

\hr

\T{5.6 \quad Last-minute recall}

\sbs{%
\lead{Say these without thinking}
\begin{itemize}
  \item $h_d[n]=\dfrac{\sin(\omega_c(n-\alpha))}{\pi(n-\alpha)}$, centre tap
        $\omega_c/\pi$, $\alpha=\dfrac{N-1}{2}$.
  \item $\omega_c=\dfrac{\omega_p+\omega_s}{2}$, and $\Delta\omega=|\omega_s-\omega_p|$.
  \item Window from \textbf{attenuation}: 21 / 25 / 44 / 53 / 74 dB.
  \item Length from \textbf{transition}: $1.8$ / $6.1$ / $6.2$ / $6.6$ / $11$, all
        $\times\pi/\Delta\omega$.
  \item $A=-20\log_{10}\delta$. Kaiser:
        $\beta=0.5842(A-21)^{0.4}+0.07886(A-21)$ for $21<A\le50$, and
        $M=\dfrac{A-8}{2.285\Delta\omega}$.
  \item Gibbs overshoot is \textbf{9 \%} and \mk{does not depend on $N$}.
  \item Alternation theorem: \textbf{$r+2$} equal-magnitude alternating extrema.
  \item $\omega=2\pi f/f_s$ for a spec in Hz.
\end{itemize}}{%
\lead{Mistakes that cost marks}
\begin{itemize}
  \item Using $\omega_p$ instead of $\omega_c$ in $h_d[n]$.
  \item \mk{Evaluating the centre tap as $0/0$} instead of $\omega_c/\pi$.
  \item Choosing the window from the transition width and the length from the
        attenuation. It is the other way round.
  \item Taking an even $N$, which puts $\alpha$ half way between samples.
  \item Saying Gibbs ripple can be fixed by a longer filter. \mk{It cannot.}
  \item Quoting $\delta_p$ when the spec gives $\delta_s$ smaller: Kaiser uses
        $\min(\delta_p,\delta_s)$.
  \item Writing $J_0$ for the Bessel function in the Kaiser window. \mk{It is $I_0$},
        the \emph{modified} one, and one paper mislabels it.
\end{itemize}}
